\documentclass[12pt]{article}
\usepackage[T1]{fontenc}
\usepackage[utf8]{inputenc}
\usepackage{lmodern}
\usepackage{textcomp}
\usepackage{enumitem}
\usepackage{geometry}
\geometry{margin=1in}
\begin{document}
\begin{center}
Answer Key - Economics - Graduate Level Assessment 16 \
\small (Answers are ordered exactly as in the question paper.)
\end{center}

\section*{Multiple Choice}
\begin{enumerate}[label=\arabic*.]
\item \textbf{Answer:} A
\\ \textit{Options:} A) 1 and 0.5; B) 2 and 0.5; C) -1 and -0.5; D) -1 and 2; E) 3 and -2
\\ \textit{Note:} The characteristic roots of the MA process are derived from the equation formed by setting the polynomial of the lag operator equal to zero, where the given MA(2) process can be expressed as \( y_t = u_t - 3 u_{t-1} + u_{t-2} \), leading to the characteristic equation \( z^2 - 3 z + 1 = 0 \), which factors to yield roots of 1 and 0.5.
\item \textbf{Answer:} A
\\ \textit{Options:} A) Buffer stocks eliminate the need for agricultural subsidies.; B) Buffer stocks result in price fluctuation.; C) Buffer stocks ensure a constant decrease in the market price of crops.; D) Buffer stocks provide unlimited income for farmers.; E) Buffer stocks are primarily used to control the export of crops.
\\ \textit{Note:} Buffer stocks stabilize agricultural prices by absorbing excess supply during periods of surplus, thereby reducing the need for government intervention through subsidies to support farmers' incomes.
\item \textbf{Answer:} A
\\ \textit{Options:} A) The Dickey-Fuller test; B) The Jarque-Bera test; C) White's test; D) The Augmented Dickey-Fuller test; E) The Durbin Watson test
\\ \textit{Note:} The Dickey-Fuller test is designed to test for the presence of unit roots in time series data, which is essential for detecting autocorrelation up to third order, as it assesses whether a time series variable is stationary or has a stochastic trend.
\item \textbf{Answer:} A
\\ \textit{Options:} A) A method that minimizes the time taken to produce a product; B) A method that uses more of all inputs to increase output; C) A method that produces the highest quantity of output; D) A method that relies solely on the latest technology; E) A method that achieves the fastest production speed regardless of resource usage
\\ \textit{Note:} The technically efficient method of production refers to the optimal use of resources to minimize the time and inputs required to produce a given output, thus maximizing productivity.
\item \textbf{Answer:} A
\\ \textit{Options:} A) the marginal propensity to save is .80.; B) the marginal propensity to save is .90.; C) the marginal propensity to save is .20.; D) the marginal propensity to consume is .20.; E) the marginal propensity to save is .50.
\\ \textit{Note:} The marginal propensity to save (MPS) is calculated as the change in saving divided by the change in disposable income, so with a $10 increase in saving for every $100 increase in income, the MPS is 10/100 = 0.10, and since the marginal propensity to consume (MPC) is 0.90 (as consumption rises by $90 when income rises by $100), the MPS, being 1 minus the MPC, is 0.80.
\end{enumerate}

\section*{True / False}
\begin{enumerate}[label=\arabic*.]
\setcounter{enumi}{5}
\item \textbf{Answer:} False
\\ \textit{Options:} A) True; B) False
\\ \textit{Note:} Stagflation, characterized by stagnant economic growth and high inflation, typically results from supply shocks or leftward shifts in the aggregate supply curve, rather than a rightward shift in aggregate demand, which would usually lead to economic growth and lower unemployment.
\item \textbf{Answer:} False
\\ \textit{Options:} A) True; B) False
\\ \textit{Note:} The statement is false because while inventors create new products, innovators play a crucial role in developing, marketing, and scaling these inventions for commercial success, highlighting the collaborative nature of the innovation process.
\item \textbf{Answer:} True
\\ \textit{Options:} A) True; B) False
\\ \textit{Note:} Disposable income is defined as the portion of an individual's income that is available for spending and saving after personal taxes have been deducted, making the statement true.
\item \textbf{Answer:} False
\\ \textit{Options:} A) True; B) False
\\ \textit{Note:} Operating at the natural rate of unemployment signifies that the economy is functioning at its potential output with normal levels of frictional and structural unemployment, rather than indicating a recession characterized by high unemployment and low economic activity.
\item \textbf{Answer:} False
\\ \textit{Options:} A) True; B) False
\\ \textit{Note:} A decrease in the price of aluminum, a key input in bicycle production, would typically lower production costs, incentivizing producers to increase supply rather than decrease it.
\end{enumerate}

\section*{Long Form Response}
\begin{enumerate}[label=\arabic*.]
\setcounter{enumi}{10}
\item \textbf{Answer:} In a capitalist market economy, the allocation of economic resources is significantly influenced by the decisions made by the central bank, particularly through its monetary policy. By adjusting interest rates and controlling money supply, the central bank can either stimulate or restrain economic activity, directly affecting investment decisions and consumption patterns. For instance, lower interest rates can encourage borrowing and spending, leading to increased demand for goods and services, thus promoting growth. However, if these actions lead to excessive inflation or asset bubbles, they can undermine market efficiency by distorting resource allocation. Ultimately, the central bank's policies play a critical role in balancing growth and stability, impacting the overall health of the economy.
\\ \textit{Note:} correct because it accurately describes how the central bank's monetary policy, through interest rate adjustments and money supply control, directly influences investment and consumption decisions in a capitalist market economy, thereby impacting both market efficiency and economic growth.
\item \textbf{Answer:} The relationship between patent issuance and monopoly power in markets is complex, as patents do not inherently create monopoly power. Instead, patents grant inventors exclusive rights to their innovations for a limited time, which can incentivize competition by encouraging firms to invest in research and development. While patents can lead to temporary monopolistic positions, they are designed to stimulate innovation by allowing inventors to recoup their costs before competitors enter the market. Consequently, while patents can restrict competition in the short term, they ultimately promote a dynamic environment where innovation flourishes, benefiting the economy as a whole. Thus, the issuance of patents can be seen as a double-edged sword; it provides temporary market control while fostering long-term competition and technological advancement.
\\ \textit{Note:} The answer correctly highlights that patents provide temporary exclusivity to inventors, fostering innovation and competition by enabling firms to recover their research and development investments, thereby illustrating the nuanced relationship between patent issuance and monopoly power in markets.
\end{enumerate}

\vfill
\noindent\textit{End of Answer Key}
\end{document}